% Preámbulo común para las soluciones de los prácticos de Programación 3
\documentclass[11pt,a4paper]{article}

\usepackage[utf8]{inputenc}
\usepackage[T1]{fontenc}
\usepackage{lmodern}
\usepackage[spanish,es-noshorthands]{babel}
\usepackage[margin=2.3cm]{geometry}
\usepackage{amsmath,amssymb,amsthm}
\usepackage{enumitem}
\usepackage{xcolor}
\usepackage{tikz}
\usetikzlibrary{arrows.meta,positioning,shapes,calc}
\usepackage[noend]{algpseudocode}
\usepackage{algorithm}
\usepackage{booktabs}
\usepackage{tcolorbox}
\usepackage[hidelinks]{hyperref}

\floatname{algorithm}{Algoritmo}
\renewcommand{\algorithmicrequire}{\textbf{Entrada:}}
\renewcommand{\algorithmicensure}{\textbf{Salida:}}
\algrenewcommand\algorithmicwhile{\textbf{mientras}}
\algrenewcommand\algorithmicdo{\textbf{hacer}}
\algrenewcommand\algorithmicfor{\textbf{para}}
\algrenewcommand\algorithmicforall{\textbf{para todo}}
\algrenewcommand\algorithmicif{\textbf{si}}
\algrenewcommand\algorithmicthen{\textbf{entonces}}
\algrenewcommand\algorithmicelse{\textbf{sino}}
\algrenewcommand\algorithmicreturn{\textbf{retornar}}
\algrenewcommand\algorithmicfunction{\textbf{función}}
\algrenewcommand\algorithmicprocedure{\textbf{procedimiento}}

\theoremstyle{plain}
\newtheorem{lema}{Lema}[section]
\newtheorem{teorema}[lema]{Teorema}
\newtheorem{proposicion}[lema]{Proposición}
\theoremstyle{definition}
\newtheorem{definicion}[lema]{Definición}
\theoremstyle{remark}
\newtheorem*{observacion}{Observación}

% Caja para el enunciado de cada ejercicio
\newtcolorbox{enunciado}{colback=gray!6,colframe=gray!50,boxrule=0.4pt,arc=2pt,left=6pt,right=6pt,top=4pt,bottom=4pt,fontupper=\small}
% Nota que vincula el ejercicio con parciales/exámenes
\newcommand{\examref}[1]{\par\noindent{\small\color{blue!50!black}\textbf{Relación con parciales/exámenes:} #1}\par\medskip}

\newcommand{\NP}{\mathsf{NP}}
\newcommand{\Pclass}{\mathsf{P}}
\newcommand{\Oh}{\mathcal{O}}

\setlength{\parskip}{0.4em}
\setlength{\parindent}{0pt}


\newcommand{\OPT}{\mathrm{OPT}}
\newcommand{\True}{\textsc{verdadero}}
\newcommand{\False}{\textsc{falso}}

\title{Práctico 9 --- Programación Dinámica II}
\author{Programación 3 --- FING, UdelaR}
\date{}

\begin{document}
\maketitle

%======================================================================
\section{Ejercicio 1}

\begin{enunciado}
\textbf{(Kleinberg \& Tardos, Ex.~6.19).} Usted es consultor de un grupo secreto de personas cuyo trabajo
consiste en monitorear y analizar señales eléctricas de barcos en la costa Atlántica. Ellos quieren un
algoritmo eficiente para el siguiente problema: desenredar un entrelazado de dos señales conocidas.

Sean $s$, $x$, $y$ cadenas binarias del mismo largo. Decimos que $s$ es un \emph{entrelazado} de $x$ e $y$
si sus símbolos se pueden particionar en dos secuencias (no necesariamente contiguas), $s'$ y $s''$, tales
que $s'$ es un prefijo de $x$ y $s''$ es un prefijo de $y$. Por ejemplo, si $x = 10110110110$ e
$y = 00100010011$, entonces $s = 10010000111$ es un entrelazado de $x$ e $y$, ya que la concatenación de los
caracteres subrayados es un prefijo de $y$, y la concatenación de los restantes es un prefijo de $x$.

\begin{enumerate}[label=(\alph*)]
\item Presente un algoritmo de programación dinámica que, dadas tres cadenas de largo $n$, $x$, $y$, $s$,
  determina si $s$ es un entrelazado de $x$ e $y$.

  \emph{Sugerencia:} Observe que si $s$ es un entrelazado de $x$ e $y$, cada carácter $c$ de $s$ proviene o
  bien de $x$ o de $y$. Defina una recurrencia para una función $M$ tal que, para cada par de naturales
  $i, j$, $i + j \le n$, el valor de $M(i, j)$ determina si $s_1 \dots s_{i+j}$ es un entrelazado de
  $x_1 \dots x_i$ y $y_1 \dots y_j$.
\item Presente un algoritmo que, si efectivamente $s$ es un entrelazado de $x$ e $y$, devuelve un conjunto
  de posiciones de $s$ que provienen de $x$ (el complemento necesariamente proviene de $y$).
\end{enumerate}
\end{enunciado}

\examref{No aparece tal cual en los parciales/exámenes revisados. La técnica (subproblemas indexados por un
par de prefijos y tabla bidimensional) es la misma del alineamiento de secuencias / distancia de edición,
que sí aparece, por ejemplo, en el Examen de febrero de 2019 (problema de distancia de edición con
$\OPT(i,j)$ sobre los prefijos $x_1\dots x_i$ e $y_1\dots y_j$).}

\subsection*{Parte (a)}

\paragraph{Subproblemas.} Para $i,j \ge 0$ con $i+j\le n$ definimos
\[
M(i,j) = \begin{cases}
\True & \text{si } s_1\dots s_{i+j} \text{ es un entrelazado de } x_1\dots x_i \text{ y } y_1\dots y_j,\\
\False & \text{en otro caso.}
\end{cases}
\]
Es decir, $M(i,j)$ es verdadero si y sólo si las posiciones $1,\dots,i+j$ de $s$ se pueden particionar en
dos subsecuencias, una igual a $x_1\dots x_i$ y la otra igual a $y_1\dots y_j$ (ambas leídas en el orden de
$s$).

Como $s'$ y $s''$ particionan las $n$ posiciones de $s$, si $s'$ tiene largo $i$ entonces $s''$ tiene largo
$n-i$, y por lo tanto
\[
s \text{ es entrelazado de } x \text{ e } y \iff \bigvee_{i=0}^{n} M(i,\,n-i).
\]

\paragraph{Recurrencia.}
\begin{itemize}
\item Caso base: $M(0,0)=\True$ (la cadena vacía es entrelazado de dos cadenas vacías).
\item Para $i+j\ge 1$:
\[
M(i,j) \;=\; \bigl(i>0 \;\wedge\; s_{i+j}=x_i \;\wedge\; M(i-1,j)\bigr)
\;\vee\;
\bigl(j>0 \;\wedge\; s_{i+j}=y_j \;\wedge\; M(i,j-1)\bigr).
\]
\end{itemize}

\paragraph{Justificación.} Supongamos $i+j\ge 1$ y que $s_1\dots s_{i+j}$ es entrelazado de
$x_1\dots x_i$ e $y_1\dots y_j$. El último símbolo $s_{i+j}$ pertenece a una de las dos subsecuencias, y
como es el último de $s_1\dots s_{i+j}$, es el último símbolo de esa subsecuencia:
\begin{itemize}
\item si pertenece a la que forma $x_1\dots x_i$, entonces $i>0$, $s_{i+j}=x_i$, y quitándolo queda una
  partición de $s_1\dots s_{i+j-1}$ en $x_1\dots x_{i-1}$ e $y_1\dots y_j$, o sea $M(i-1,j)$;
\item si pertenece a la que forma $y_1\dots y_j$, análogamente $j>0$, $s_{i+j}=y_j$ y $M(i,j-1)$.
\end{itemize}
Recíprocamente, si por ejemplo $i>0$, $s_{i+j}=x_i$ y $M(i-1,j)$, tomando la partición de
$s_1\dots s_{i+j-1}$ que testimonia $M(i-1,j)$ y agregando la posición $i+j$ a la subsecuencia de $x$ se
obtiene una partición válida de $s_1\dots s_{i+j}$; lo mismo para el otro caso. Por lo tanto la
recurrencia es correcta.

\paragraph{Algoritmo iterativo.} Cada $M(i,j)$ depende sólo de $M(i-1,j)$ y $M(i,j-1)$, por lo que alcanza
con recorrer la tabla por filas crecientes en $i$ y, dentro de cada fila, por $j$ creciente (o, de manera
equivalente, por valores crecientes de $i+j$).

\begin{algorithm}[H]
\caption{Entrelazado$(x,y,s)$}
\begin{algorithmic}[1]
\Require cadenas $x,y,s$ de largo $n$
\Ensure \True{} si $s$ es entrelazado de $x$ e $y$ (y la tabla $M$)
\State $M[0,0] \gets \True$
\For{$i \gets 0$ \textbf{to} $n$}
  \For{$j \gets 0$ \textbf{to} $n-i$}
    \If{$(i,j)\neq(0,0)$}
      \State $M[i,j] \gets (i>0 \wedge s_{i+j}=x_i \wedge M[i-1,j]) \vee (j>0 \wedge s_{i+j}=y_j \wedge M[i,j-1])$
    \EndIf
  \EndFor
\EndFor
\For{$i \gets 0$ \textbf{to} $n$}
  \If{$M[i,n-i]$} \Return \True \EndIf
\EndFor
\State \Return \False
\end{algorithmic}
\end{algorithm}

\paragraph{Orden.} La tabla tiene $\sum_{i=0}^n (n-i+1) = \frac{(n+1)(n+2)}{2}$ celdas y cada una se calcula
en $\Oh(1)$; la recorrida final es $\Oh(n)$. El tiempo total es $\Oh(n^2)$ y el espacio $\Oh(n^2)$ (si sólo
se quiere la respuesta sí/no basta con guardar la fila anterior: espacio $\Oh(n)$).

\subsection*{Parte (b)}

Si $s$ es entrelazado, existe $i_0$ con $M(i_0,n-i_0)=\True$. Partimos de $(i,j)=(i_0,n-i_0)$ y
retrocedemos usando la recurrencia: en cada paso, $M(i,j)=\True$ con $i+j\ge1$, así que al menos uno de los
dos disyuntos es verdadero. Si es el primero, la posición $i+j$ de $s$ proviene de $x$ y pasamos a
$(i-1,j)$; si no, proviene de $y$ y pasamos a $(i,j-1)$. Se mantiene el invariante $M(i,j)=\True$, por lo
que el proceso termina en $(0,0)$ tras exactamente $n$ pasos y, por la justificación de la recurrencia, las
posiciones marcadas forman una partición válida: las asignadas a $x$ forman $x_1\dots x_{i_0}$ y las
restantes $y_1\dots y_{n-i_0}$.

\begin{algorithm}[H]
\caption{PosicionesDeX$(x,y,s)$}
\begin{algorithmic}[1]
\State calcular la tabla $M$ con el Algoritmo Entrelazado
\State sea $i_0$ tal que $M[i_0,n-i_0]=\True$ \Comment{si no existe, $s$ no es entrelazado}
\State $P \gets \emptyset$; \ $i\gets i_0$; \ $j \gets n-i_0$
\While{$i+j>0$}
  \If{$i>0 \wedge s_{i+j}=x_i \wedge M[i-1,j]$}
    \State $P \gets P\cup\{i+j\}$; \ $i\gets i-1$
  \Else \Comment{necesariamente $j>0$, $s_{i+j}=y_j$ y $M[i,j-1]$}
    \State $j\gets j-1$
  \EndIf
\EndWhile
\State \Return $P$
\end{algorithmic}
\end{algorithm}

La reconstrucción hace $n$ iteraciones de costo $\Oh(1)$, así que el total sigue siendo $\Oh(n^2)$ (dominado
por el cálculo de la tabla).

\paragraph{Ejemplo.} Para $x = 10110110110$, $y = 00100010011$, $s = 10010000111$ ($n=11$), la tabla queda
como sigue (T $=\True$, $\cdot=\False$; fila $i$, columna $j$, sólo $i+j\le 11$):

\begin{center}
\small
\setlength{\tabcolsep}{4pt}
\begin{tabular}{c|cccccccccccc}
$i\backslash j$ & 0 & 1 & 2 & 3 & 4 & 5 & 6 & 7 & 8 & 9 & 10 & 11\\ \hline
0  & T & $\cdot$ & $\cdot$ & $\cdot$ & $\cdot$ & $\cdot$ & $\cdot$ & $\cdot$ & $\cdot$ & $\cdot$ & $\cdot$ & $\cdot$\\
1  & T & T & T & T & T & T & T & $\cdot$ & $\cdot$ & $\cdot$ & $\cdot$ & \\
2  & T & T & $\cdot$ & T & T & T & T & T & $\cdot$ & $\cdot$ & & \\
3  & $\cdot$ & T & T & $\cdot$ & $\cdot$ & $\cdot$ & T & T & $\cdot$ & & & \\
4  & $\cdot$ & $\cdot$ & $\cdot$ & $\cdot$ & $\cdot$ & $\cdot$ & T & \textbf{T} & & & & \\
5  & $\cdot$ & $\cdot$ & $\cdot$ & $\cdot$ & $\cdot$ & $\cdot$ & $\cdot$ & & & & & \\
6  & $\cdot$ & $\cdot$ & $\cdot$ & $\cdot$ & $\cdot$ & $\cdot$ & & & & & & \\
7  & $\cdot$ & $\cdot$ & $\cdot$ & $\cdot$ & $\cdot$ & & & & & & & \\
8  & $\cdot$ & $\cdot$ & $\cdot$ & $\cdot$ & & & & & & & & \\
9  & $\cdot$ & $\cdot$ & $\cdot$ & & & & & & & & & \\
10 & $\cdot$ & $\cdot$ & & & & & & & & & & \\
11 & $\cdot$ & & & & & & & & & & & \\
\end{tabular}
\end{center}

En la antidiagonal $i+j=11$ la única celda verdadera es $M(4,7)$, así que $s$ es entrelazado (con 4
símbolos tomados de $x$ y 7 de $y$). Retrocediendo desde $(4,7)$ el algoritmo devuelve
$P=\{1,8,10,11\}$: en efecto $s_1s_8s_{10}s_{11} = 1011 = x_1\dots x_4$ y las posiciones restantes
$2,\dots,7,9$ dan $0010001 = y_1\dots y_7$.

%======================================================================
\section{Ejercicio 2}

\begin{enunciado}
\textbf{(Kleinberg \& Tardos, Ex.~6.21).} Una compañía de inversiones está analizando el valor de las
acciones de una empresa que cotiza en bolsa a lo largo de $n$ días consecutivos. Los días considerados están
numerados de $1$ a $n$; para cada día $i$, $1\le i\le n$, cada acción tiene un precio $p(i)$.

Para algunos valores, posiblemente grandes, de un parámetro $k$, quieren estudiar lo que llaman
\emph{estrategias de $k$-disparos}. Una estrategia de $k$-disparos es una colección de $m$ pares de días
$(c_1,v_1),\dots,(c_m,v_m)$ donde $0\le m\le k$ y
\[ 1\le c_1<v_1<c_2<v_2<\dots<c_m<v_m\le n. \]
Podemos interpretar esta secuencia como un conjunto de $m$ intervalos disjuntos, durante los que los
inversores compran 1000 acciones del stock (en el día $c_i$) y luego las venden (en el día $v_i$). Definimos
entonces la ganancia de una estrategia de $k$-disparos como
\[ \mathrm{ganancia}(E) = 1000\cdot\sum_{i=1}^{m}\bigl(p(v_i)-p(c_i)\bigr), \]
donde $E$ es una estrategia de $k$-disparos para un stock de una acción dada.
\begin{enumerate}[label=(\alph*)]
\item Diseñe un algoritmo eficiente que, dada una secuencia de precios, determina una estrategia de
  $k$-disparos con la mayor ganancia posible para una acción dada. Asuma que siempre se cuenta con stock
  suficiente para comprar 1000 acciones en un día $i$. El tiempo de ejecución de su algoritmo debe ser
  polinomial en $n$ y $k$, es decir que no debe contener a $k$ en el exponente.
\item Analice el tiempo de ejecución de su algoritmo en función de $n$ y $k$.
\end{enumerate}
\emph{Sugerencia:} Para cada $i$, $0\le i\le n$, calcule la mayor ganancia que es posible obtener mediante
una estrategia de $d$-disparos, con $0\le d\le k$, a lo largo de los días $1$ hasta $i$. Para ello evalúe la
conveniencia de vender las acciones del último disparo en el día $i$ o no hacerlo.
\end{enunciado}

\examref{No encontramos este problema en los parciales/exámenes revisados. Es un problema de
programación dinámica en dos parámetros del estilo de mochila (\emph{subset sum}): un parámetro recorre
los días y el otro la cantidad de disparos todavía disponibles.}

Como el factor $1000$ multiplica a todas las estrategias, maximizar $\mathrm{ganancia}(E)$ equivale a maximizar
$g(E)=\sum_{t=1}^m (p(v_t)-p(c_t))$; trabajamos con $g$ y al final multiplicamos por 1000.

\subsection*{Parte (a)}

\paragraph{Subproblemas.} Para $0\le d\le k$ y $0\le i\le n$ sea
\[
\OPT(d,i) = \max\{\, g(E) : E \text{ estrategia de } d\text{-disparos con } v_m\le i\,\},
\]
o sea, la mayor ganancia con a lo sumo $d$ pares de compra-venta usando sólo los días $1,\dots,i$. El valor
buscado es $1000\cdot\OPT(k,n)$. Notar que la estrategia vacía ($m=0$) siempre es válida, así que
$\OPT(d,i)\ge 0$.

\paragraph{Recurrencia.}
\begin{itemize}
\item $\OPT(0,i)=0$ para todo $i$ (sin disparos sólo está la estrategia vacía).
\item $\OPT(d,0)=0$ para todo $d$ (no hay días).
\item Para $d\ge1$, $i\ge1$:
\[
\OPT(d,i) = \max\Bigl\{\; \OPT(d,i-1),\;\; \max_{1\le c<i}\bigl(\OPT(d-1,c-1) + p(i)-p(c)\bigr) \Bigr\},
\]
donde el máximo interno vale $-\infty$ si $i=1$ (no hay día de compra posible).
\end{itemize}

\paragraph{Justificación.} Sea $E$ una estrategia óptima para $(d,i)$ con $d,i\ge 1$. Hay dos casos,
según si se vende en el día $i$ o no:
\begin{itemize}
\item \emph{No hay venta en el día $i$} (en particular si $m=0$): entonces $v_m\le i-1$ y $E$ es una
  estrategia de $d$-disparos sobre los días $1,\dots,i-1$, por lo que $g(E)\le \OPT(d,i-1)$.
\item \emph{El último disparo vende en el día $i$}: $v_m=i$ y sea $c=c_m$, con $1\le c<i$. Los $m-1\le d-1$
  pares restantes cumplen $v_{m-1}<c$, o sea, forman una estrategia de $(d-1)$-disparos sobre los días
  $1,\dots,c-1$. Luego $g(E) = g(E\setminus\{(c,i)\}) + p(i)-p(c) \le \OPT(d-1,c-1)+p(i)-p(c)$.
\end{itemize}
Recíprocamente, cada término del máximo es la ganancia de una estrategia válida para $(d,i)$: una óptima
para $(d,i-1)$ es válida para $(d,i)$, y a una óptima para $(d-1,c-1)$ (todas sus ventas son $\le c-1<c$) se le
puede agregar el par $(c,i)$ respetando el orden estricto y sin superar los $d$ pares. Por lo tanto
$\OPT(d,i)$ es exactamente el máximo de la recurrencia.

\paragraph{Mejora a $\Oh(1)$ por celda.} Calcular directamente el máximo interno cuesta $\Oh(i)$ por celda.
Definimos
\[
B(d,i)=\max_{1\le c<i}\bigl(\OPT(d-1,c-1)-p(c)\bigr) \qquad(B(d,1)=-\infty),
\]
que es el mejor ``saldo'' tras haber comprado por última vez antes del día $i$. Se cumple
$B(d,i+1)=\max\{B(d,i),\ \OPT(d-1,i-1)-p(i)\}$, y la recurrencia queda
\[
\OPT(d,i)=\max\{\OPT(d,i-1),\ B(d,i)+p(i)\}.
\]
Guardamos también en $C[d,i]$ el día de compra $c$ que realiza $B(d,i)$, para reconstruir la estrategia.

\paragraph{Algoritmo iterativo.} $\OPT(d,i)$ depende de $\OPT(d,i-1)$ (misma fila) y de valores de la fila
$d-1$; por lo tanto llenamos la tabla por filas $d=0,1,\dots,k$ y en cada fila por $i$ creciente.

\begin{algorithm}[H]
\caption{KDisparos$(p(1..n),k)$}
\begin{algorithmic}[1]
\State $k \gets \min(k,\lfloor n/2\rfloor)$ \Comment{nunca se usan más de $\lfloor n/2\rfloor$ pares}
\For{$i\gets 0$ \textbf{to} $n$} $\OPT[0,i]\gets 0$ \EndFor
\For{$d\gets 1$ \textbf{to} $k$}
  \State $\OPT[d,0]\gets 0$; \ $best\gets -\infty$; \ $bc \gets \textsc{nil}$
  \For{$i \gets 1$ \textbf{to} $n$}
    \If{$i\ge 2$ \textbf{and} $\OPT[d-1,i-2]-p(i-1) > best$} \Comment{nuevo candidato $c=i-1$}
      \State $best \gets \OPT[d-1,i-2]-p(i-1)$; \ $bc\gets i-1$
    \EndIf
    \State $C[d,i]\gets bc$ \Comment{ahora $best = B(d,i)$}
    \If{$best + p(i) > \OPT[d,i-1]$}
      \State $\OPT[d,i]\gets best+p(i)$; \ $V[d,i]\gets \True$ \Comment{se vende en el día $i$}
    \Else
      \State $\OPT[d,i]\gets \OPT[d,i-1]$; \ $V[d,i]\gets\False$
    \EndIf
  \EndFor
\EndFor
\State \Comment{Reconstrucción}
\State $E\gets\emptyset$; \ $d\gets k$; \ $i\gets n$
\While{$d>0$ \textbf{and} $i>0$}
  \If{$V[d,i]$}
    \State $c\gets C[d,i]$; \ $E\gets E\cup\{(c,i)\}$; \ $i\gets c-1$; \ $d\gets d-1$
  \Else
    \State $i\gets i-1$
  \EndIf
\EndWhile
\State \Return $E$ \ (con ganancia $1000\cdot\OPT[k,n]$)
\end{algorithmic}
\end{algorithm}

Si $k=0$ el algoritmo devuelve la estrategia vacía, con ganancia $0$. La reconstrucción es correcta porque
en cada paso sigue la opción que realizó el máximo: si $V[d,i]$, el valor $\OPT(d,i)$ se obtuvo como
$\OPT(d-1,c-1)+p(i)-p(c)$ con $c=C[d,i]$, y continuamos con el subproblema $(d-1,c-1)$; si no, con $(d,i-1)$.
Los pares se agregan en orden decreciente de días y cumplen $c<i$ y la venta siguiente es $\le c-1$, de modo
que $E$ respeta el orden estricto requerido.

\paragraph{Ejemplo.} Precios $p=(5,2,6,1,4,3,8)$ ($n=7$), $k=3$. La tabla $\OPT(d,i)$ (en miles) queda:

\begin{center}
\begin{tabular}{c|cccccccc}
$d\backslash i$ & 0 & 1 & 2 & 3 & 4 & 5 & 6 & 7\\ \hline
0 & 0 & 0 & 0 & 0 & 0 & 0 & 0 & 0\\
1 & 0 & 0 & 0 & 4 & 4 & 4 & 4 & 7\\
2 & 0 & 0 & 0 & 4 & 4 & 7 & 7 & 11\\
3 & 0 & 0 & 0 & 4 & 4 & 7 & 7 & 12\\
\end{tabular}
\end{center}

Por ejemplo, $\OPT(2,7)=\OPT(1,3)+p(7)-p(4)=4+8-1=11$. Con $k=3$ la reconstrucción da
$E=\{(2,3),(4,5),(6,7)\}$, con ganancia $1000\cdot(4+3+5)=12000$.

\subsection*{Parte (b)}

La tabla tiene $(k+1)(n+1)$ celdas, y gracias a la variable $best$ cada una se calcula en $\Oh(1)$; la
reconstrucción hace a lo sumo $n$ iteraciones. El tiempo total es $\Oh(nk)$ y el espacio $\Oh(nk)$ (para
reconstruir hacen falta $V$ y $C$; si sólo se quiere el valor alcanza con dos filas, $\Oh(n)$).

Si se usara la recurrencia sin la variable auxiliar $B$, cada celda costaría $\Oh(n)$ y el total sería
$\Oh(n^2k)$, que también es polinomial en $n$ y $k$. Además, como una estrategia tiene a lo sumo
$\lfloor n/2\rfloor$ pares (usa $2m$ días distintos), el algoritmo reemplaza $k$ por
$\min(k,\lfloor n/2\rfloor)$, de modo que el tiempo es en realidad $\Oh(n\cdot\min(k,n))\subseteq\Oh(n^2)$:
aunque $k$ sea ``muy grande'', el tiempo no depende de él.

%======================================================================
\section{Ejercicio 3}

\begin{enunciado}
\textbf{(Kleinberg \& Tardos, Ex.~6.22).} Consideramos un grafo dirigido $G=(V,E)$ con aristas ponderadas.
Los costos en las aristas pueden ser negativos o positivos, pero todo ciclo en $G$ tiene costo positivo.
Para medir la robustez en la conectividad de $G$ queremos contar la cantidad de caminos de costo mínimo entre
dos vértices. El objetivo de este ejercicio es construir un algoritmo que resuelva el problema en tiempo
polinomial utilizando programación dinámica. El algoritmo toma como entrada un par de vértices de $G$, $s$ y
$t$, y da como resultado un número natural que es la cantidad de caminos de costo mínimo de $s$ a $t$.

Definimos $\OPT(i,v)$ como el costo de los caminos de costo mínimo de $s$ a $v$ compuestos por exactamente
$i$ aristas. Definimos también $N(i,v)$ como la cantidad de esos caminos de costo mínimo.
\begin{enumerate}[label=(\alph*)]
\item Teniendo en cuenta que un camino desde $s$ hasta $v$ de $i$ aristas, $i>0$, contiene un camino desde
  $s$ hasta un nodo adyacente a $v$ de $i-1$ aristas, escriba una relación de recurrencia que defina
  $\OPT(i,v)$ en función de los valores de $\OPT(i-1,w)$, $w\in V$.
\item Asumiendo conocidos los valores de $\OPT$, escriba una relación de recurrencia para $N(i,v)$, $i>0$,
  $v\in V$, en función de los valores $N(i-1,w)$, $w\in V$, y la función $\OPT$.
\item Observe que el costo de los caminos de costo mínimo de $s$ a $t$ viene dado por
  \[ \OPT(t)=\min_i\{\OPT(i,t)\},\qquad \OPT(s)=0. \]
  Defina de forma análoga cómo se calcula la cantidad de caminos de costo mínimo, $N(t)$, en función de los
  $N(i,t)$.
\item Escriba un algoritmo que devuelve la cantidad de caminos de costo mínimo. Para ello inicialice los
  valores de $\OPT(0,v)$ y $N(0,v)$ como crea conveniente. Luego compute las recurrencias de las partes
  anteriores y devuelva el resultado.
\end{enumerate}
\end{enunciado}

\examref{Es una generalización del algoritmo de Bellman--Ford (sección 6.8 del libro; aquí $\OPT(i,v)$ se
mide desde el origen $s$ y con \emph{exactamente} $i$ aristas). La idea de contar caminos mínimos sumando las
cantidades de los predecesores que realizan el mínimo aparece, para grafos sin pesos (con BFS), en el Examen
de julio de 2019, Ej.~2 (cantidad de historias de largo mínimo, $N(v)$ en función de los $N(u)$ del nivel
anterior) y en el Examen de julio de 2014 (cantidad de caminos de longitud mínima desde $u$). En el Examen
de febrero de 2026 se pide además la tabla de $\OPT(i,v)$ de Bellman--Ford.}

\paragraph{Notación y una observación previa.} Llamamos $c(w,v)$ al costo de la arista $(w,v)$, $n=|V|$,
$m=|E|$. En las definiciones de $\OPT(i,v)$ y $N(i,v)$ entendemos ``camino'' en sentido amplio
(\emph{recorrido}, que puede repetir vértices): con $i$ fijo, el recorrido más barato de exactamente $i$
aristas puede no ser simple, y es justamente para recorridos que vale la subestructura óptima de la parte (a)
(el prefijo de un recorrido es un recorrido cualquiera, sin restricción de no pasar por $v$). Si no existe
ningún recorrido de $s$ a $v$ con $i$ aristas, $\OPT(i,v)=\infty$ y $N(i,v)=0$. Para el problema original
(caminos simples) usamos el siguiente lema.

\begin{lema}\label{lema:simple}
Si todo ciclo de $G$ tiene costo positivo, todo recorrido de costo mínimo entre los recorridos de $s$ a $t$
es un camino simple. En consecuencia el costo mínimo de un recorrido de $s$ a $t$ existe, coincide con el
costo mínimo de un camino simple, y se alcanza con a lo sumo $n-1$ aristas.
\end{lema}
\begin{proof}
Si un recorrido $P$ de $s$ a $t$ repite un vértice, contiene un ciclo $Z$ (el tramo entre dos apariciones
del mismo vértice); quitándolo se obtiene otro recorrido de $s$ a $t$ de costo $c(P)-c(Z)<c(P)$. Repitiendo,
todo recorrido tiene costo mayor o igual que el de algún camino simple, y si no es simple tiene costo
estrictamente mayor. Como hay finitos caminos simples, el mínimo existe, sólo lo alcanzan caminos simples, y
éstos tienen a lo sumo $n-1$ aristas.
\end{proof}

\subsection*{Parte (a)}

\begin{itemize}
\item Caso base: $\OPT(0,s)=0$ y $\OPT(0,v)=\infty$ para $v\neq s$ (el único recorrido con 0 aristas que
  sale de $s$ es el trivial, que termina en $s$).
\item Para $i\ge1$:
\[
\OPT(i,v) = \min_{w\,:\,(w,v)\in E}\bigl\{\OPT(i-1,w)+c(w,v)\bigr\},
\]
con $\min\emptyset=\infty$ e $\infty+c=\infty$.
\end{itemize}

\emph{Justificación.} Un recorrido $P$ de $s$ a $v$ con $i\ge1$ aristas termina con una arista $(w,v)$ y su
prefijo $P'$ es un recorrido de $s$ a $w$ con $i-1$ aristas; así $c(P)=c(P')+c(w,v)\ge
\OPT(i-1,w)+c(w,v)$. Recíprocamente, para cada $(w,v)\in E$, agregando la arista $(w,v)$ a un recorrido
óptimo de $i-1$ aristas hasta $w$ se obtiene un recorrido de $i$ aristas hasta $v$ de costo
$\OPT(i-1,w)+c(w,v)$. Luego el mínimo coincide.

\subsection*{Parte (b)}

\begin{itemize}
\item Caso base: $N(0,s)=1$ y $N(0,v)=0$ para $v\neq s$.
\item Para $i\ge1$: si $\OPT(i,v)=\infty$ entonces $N(i,v)=0$; si no,
\[
N(i,v) = \sum_{\substack{w\,:\,(w,v)\in E\\ \OPT(i-1,w)+c(w,v)=\OPT(i,v)}} N(i-1,w).
\]
\end{itemize}

\emph{Justificación.} Sea $P$ un recorrido de $s$ a $v$ de $i$ aristas de costo $\OPT(i,v)<\infty$, con
última arista $(w,v)$ y prefijo $P'$. Por la parte (a), $\OPT(i,v)\le\OPT(i-1,w)+c(w,v)$, de donde
$c(P')=\OPT(i,v)-c(w,v)\le\OPT(i-1,w)$; como también $c(P')\ge\OPT(i-1,w)$, resulta $c(P')=\OPT(i-1,w)$
y $\OPT(i-1,w)+c(w,v)=\OPT(i,v)$. O sea: $P'$ es un recorrido óptimo de $i-1$ aristas hasta un predecesor
$w$ que realiza el mínimo. Recíprocamente, todo recorrido óptimo hasta un tal $w$, extendido con $(w,v)$, es
un recorrido óptimo hasta $v$. Esto da una biyección entre los recorridos óptimos hasta $v$ con $i$ aristas y
los pares ($w$ que realiza el mínimo, recorrido óptimo de $i-1$ aristas hasta $w$); recorridos con distinto
$w$ o distinto prefijo son distintos. Por lo tanto la cantidad es la suma indicada.

\subsection*{Parte (c)}

Por el Lema~\ref{lema:simple} basta considerar $0\le i\le n-1$:
\[
\OPT(t)=\min_{0\le i\le n-1}\OPT(i,t),
\qquad
N(t)=\sum_{\substack{0\le i\le n-1\\ \OPT(i,t)=\OPT(t)}} N(i,t),
\]
y si $\OPT(t)=\infty$ ($t$ no es alcanzable desde $s$) entonces $N(t)=0$.

\emph{Justificación.} Todo camino de costo mínimo de $s$ a $t$ tiene una cantidad $i\le n-1$ de aristas y su
costo es $\OPT(t)$, por lo que necesariamente $\OPT(i,t)=\OPT(t)$ (no puede haber recorridos más baratos con
$i$ aristas) y está contado en $N(i,t)$. Recíprocamente, si $\OPT(i,t)=\OPT(t)$, cada recorrido contado en
$N(i,t)$ tiene costo mínimo global y, por el lema, es un camino simple; así que no contamos recorridos con
ciclos. Además los conjuntos contados para distintos $i$ son disjuntos (tienen distinta cantidad de aristas),
de modo que la suma no cuenta nada dos veces.

Para $t=s$: el camino trivial tiene costo $0$ y todo recorrido cerrado de $s$ a $s$ con $i\ge1$ aristas se
descompone en ciclos de costo positivo, así que $\OPT(s)=0$ y $N(s)=N(0,s)=1$, consistente con el enunciado.

\subsection*{Parte (d)}

Inicializamos como en las partes (a) y (b) y llenamos la tabla por $i$ creciente: la fila $i$ sólo depende
de la fila $i-1$. En cada fila recorremos las aristas una vez.

\begin{algorithm}[H]
\caption{CantCaminosMinimos$(G=(V,E),c,s,t)$}
\begin{algorithmic}[1]
\Require grafo dirigido $G$ con costos $c$ y todo ciclo de costo positivo; $s,t\in V$
\Ensure cantidad de caminos de costo mínimo de $s$ a $t$
\ForAll{$v\in V$}
  \State $\OPT[0,v]\gets\infty$; \ $N[0,v]\gets 0$
\EndFor
\State $\OPT[0,s]\gets 0$; \ $N[0,s]\gets 1$
\For{$i\gets 1$ \textbf{to} $n-1$}
  \ForAll{$v\in V$}
    \State $\OPT[i,v]\gets\infty$; \ $N[i,v]\gets 0$
  \EndFor
  \ForAll{$(w,v)\in E$} \Comment{cálculo de $\OPT(i,\cdot)$, parte (a)}
    \If{$\OPT[i-1,w]<\infty$ \textbf{and} $\OPT[i-1,w]+c(w,v)<\OPT[i,v]$}
      \State $\OPT[i,v]\gets\OPT[i-1,w]+c(w,v)$
    \EndIf
  \EndFor
  \ForAll{$(w,v)\in E$} \Comment{cálculo de $N(i,\cdot)$, parte (b)}
    \If{$\OPT[i-1,w]<\infty$ \textbf{and} $\OPT[i-1,w]+c(w,v)=\OPT[i,v]$}
      \State $N[i,v]\gets N[i,v]+N[i-1,w]$
    \EndIf
  \EndFor
\EndFor
\State $best\gets\min_{0\le i\le n-1}\OPT[i,t]$ \Comment{parte (c)}
\If{$best=\infty$} \Return $0$ \EndIf
\State $total\gets 0$
\For{$i\gets 0$ \textbf{to} $n-1$}
  \If{$\OPT[i,t]=best$} $total\gets total+N[i,t]$ \EndIf
\EndFor
\State \Return $total$
\end{algorithmic}
\end{algorithm}

\paragraph{Orden.} Hay $n-1$ iteraciones del ciclo principal, cada una de costo $\Oh(n+m)$ (inicializar la
fila y recorrer dos veces las aristas). El tiempo total es $\Oh(n(n+m))$, es decir $\Oh(nm)$ si $m\ge n$,
igual que Bellman--Ford. El espacio es $\Oh(n^2)$, o bien $\Oh(n)$ si se guardan sólo las filas $i-1$ e $i$
y se van acumulando $best$ y $total$ a medida que se calcula $\OPT[i,t]$ (reiniciando $total$ cada vez que
$best$ mejora). Observación: la cantidad de caminos puede ser exponencial en $n$, pero se representa con
$\Oh(n\log n)$ bits, así que las sumas siguen siendo de costo polinomial.

\paragraph{Ejemplo.} Sea $V=\{s,a,b,t\}$ con aristas $(s,a)$ de costo $1$, $(s,b)$ de costo $0$, $(a,b)$ de
costo $-1$, $(a,t)$ de costo $2$, $(b,t)$ de costo $3$, $(s,t)$ de costo $3$ y $(t,a)$ de costo $1$ (los
ciclos $a\,t\,a$ y $a\,b\,t\,a$ cuestan $3$). La tabla de pares $\OPT(i,v)\,/\,N(i,v)$ es:

\begin{center}
\begin{tabular}{c|cccc}
$i\backslash v$ & $s$ & $a$ & $b$ & $t$\\ \hline
0 & $0\,/\,1$ & $\infty\,/\,0$ & $\infty\,/\,0$ & $\infty\,/\,0$\\
1 & $\infty\,/\,0$ & $1\,/\,1$ & $0\,/\,1$ & $3\,/\,1$\\
2 & $\infty\,/\,0$ & $4\,/\,1$ & $0\,/\,1$ & $3\,/\,2$\\
3 & $\infty\,/\,0$ & $4\,/\,2$ & $3\,/\,1$ & $3\,/\,1$\\
\end{tabular}
\end{center}

Por ejemplo $\OPT(3,a)=4$ se alcanza por $s\,a\,t\,a$ y $s\,b\,t\,a$ (recorridos no simples, de ahí
$N(3,a)=2$), y $N(2,t)=2$ porque tanto $w=a$ ($1+2$) como $w=b$ ($0+3$) realizan el mínimo $3$. Luego
$\OPT(t)=3$ y $N(t)=N(1,t)+N(2,t)+N(3,t)=1+2+1=4$: los caminos $s\,t$, $s\,a\,t$, $s\,b\,t$ y
$s\,a\,b\,t$, todos de costo $3$.

\end{document}
